\section*{Hoofdstuk \ref{ch:b3:hydrogeology} --- Hydrogeologie}

\begin{solution}{exo:b3:hydrogeology:1}
Eerste filter: $z = 120 - 40 = \qty{80}{m}$, $h = 120 - 12.5 =
\qty{107.5}{m}$, drukhoogte $107.5 - 80 = \qty{27.5}{m}$. Tweede
filter: $z = \qty{40}{m}$, $h = \qty{105.0}{m}$, drukhoogte
\qty{65.0}{m}. De \omterm{def:b3:hydrogeology:head}{stijghoogte} is hoger bij het ondiepe filter, dus het water
stroomt naar beneden, onder een \omterm{def:b3:hydrogeology:head}{hydraulische gradiënt} $(107.5 - 105.0)/40 = \num{0.0625}$:
de peilbuizen staan in een infiltratiegebied. De drukhoogte alleen
had het tegendeel doen vermoeden.
\end{solution}

\begin{solution}{exo:b3:hydrogeology:2}
$q = K\Delta h/L = 12\times0.30/0.50 = \qty{7.2}{m/d}$. Het oppervlak is
$\pi\times0.05^2 = \qty{7.85e-3}{m^2}$, dus $Q = \qty{0.0565}{m^3/d}$, of
$56.5/24 = \qty{2.36}{L/h}$. De \omterm{def:b3:hydrogeology:seepage}{poriesnelheid} is $7.2/0.35 =
\qty{20.6}{m/d}$.
\end{solution}

\begin{solution}{exo:b3:hydrogeology:3}
$K = 2.5/86400 = \qty{2.894e-5}{m/s}$, en $k = K\mu/(\rho g) =
\num{2.894e-5}\times\num{1.306e-3}/(999.7\times9.81) =
\qty{3.85e-12}{m^2}$. Bij \qty{20}{\celsius} wordt $K$ vermenigvuldigd met
$(998.2/\num{1.002e-3})/(999.7/\num{1.306e-3}) = \num{1.30}$: $K =
\qty{3.25}{m/d}$. De verandering komt vrijwel volledig door de viscositeit.
\end{solution}

\begin{solution}{exo:b3:hydrogeology:4}
$T = 2.5\times180 = \qty{450}{m^2/d}$. De uitzetting van het water geeft
$\num{9e-7}\times180 = \num{1.62e-4}$, ongeveer \qty{80}{\percent} van $S =
\num{2e-4}$; de compressie van het kalksteen levert de rest.
\end{solution}

\begin{solution}{exo:b3:hydrogeology:5}
Volgens \cref{prop:b3:hydrogeology:thiem}, $T = Q\ln(r_2/r_1)/(2\pi(s_1 -
s_2)) = 1500\times\ln10/(2\pi\times1.53) = \qty{359}{m^2/d}$, en $K =
359/40 = \qty{9.0}{m/d}$.
\end{solution}

\begin{solution}{exo:b3:hydrogeology:6}
$Q/4\pi T = 1200/(4\pi\times300) = \qty{0.318}{m}$. Na 1 dag, $u =
300^2\times\num{5e-4}/(4\times300\times1) = \num{0.0375}$ en $W = -0.5772
+ 3.2834 + 0.0375 - 0.0004 = \num{2.743}$: $s = \qty{0.873}{m}$. Na 10
dagen, $u = \num{0.00375}$, $W = -0.5772 + 5.5860 + 0.0038 = \num{5.013}$:
$s = \qty{1.596}{m}$. Tien keer zo lang geeft minder dan twee keer de
verlaging.
\end{solution}

\begin{solution}{exo:b3:hydrogeology:7}
$u \le 0.05$ wanneer $t \ge r^2S/(4T\times0.05)$. Bij \qty{60}{m}: $3600\times
\num{2e-4}/90 = \qty{0.008}{d}$, \qty{11.5}{min}. Bij \qty{500}{m}:
\qty{0.556}{d}, \qty{800}{min}, ruim 13 uur. De tijd groeit met $r^2$:
in een test van een dag of twee moeten de waarnemingsputten voor de rechte
lijn binnen enkele honderden meter van de onttrekkingsput liggen.
\end{solution}

\begin{solution}{exo:b3:hydrogeology:8}
Halverwege staat elke put op \qty{200}{m} afstand: $u = 200^2\times\num{2e-4}/(4
\times450\times10) = \num{4.44e-4}$, $W = \num{7.142}$, $s = 0.382\times
7.142 = \qty{2.728}{m}$ per put, \qty{5.456}{m} in totaal. Bij
\qty{400}{m}, $u = \num{1.78e-3}$, $W = \num{5.757}$, $s =
\qty{2.199}{m}$. Elke put verlaagt het waterpeil in de andere met
\qty{2.2}{m}: de pomp moet het water \qty{2.2}{m} hoger opvoeren, wat
kost energie, en de verlaging die overblijft voordat de pompinlaat droogvalt, is
kleiner.
\end{solution}

\begin{solution}{exo:b3:hydrogeology:9}
Vier stroomlijnen tussen de damvoet en de ondergrond begrenzen $N_f = 5$
stroomkanalen; negen equipotentiaallijnen verdelen de \qty{10}{m} in $N_d = 10$
potentiaalvallen. $q = KHN_f/N_d = 5\times10\times5/10 = \qty{25}{m^3/d}$ per meter
van de dam, en \qty{7500}{m^3/d} onder \qty{300}{m}. De berekende
vormfactor, 0.512 in plaats van 0.5, zou \qty{25.6}{m^3/d} per meter opleveren.
\end{solution}

\begin{solution}{exo:b3:hydrogeology:10}
$k = na^2/8 = 0.30\times(\num{5e-5})^2/8 = \qty{9.4e-11}{m^2}$, en $K =
k\rho g/\mu = \num{9.375e-11}\times998.2\times9.81/\num{1.002e-3} =
\qty{9.2e-4}{m/s}$, \qty{79}{m/d}, binnen het bereik van schoon zand, 0.3 tot
\qty{150}{m/d}. Het model laat de kronkeligheid van de paden buiten beschouwing, die langer zijn
dan de rechte buizen (dat verlaagt $K$), de engten tussen de poriën,
kleiner dan de poriën (dat verlaagt), de doodlopende poriën die geen
stroming dragen (dat verlaagt), en de spreiding van de poriegroottes: omdat een buis een
stroming draagt evenredig met $a^4$, de breedste paden voeren het meeste water en
een spreiding verhoogt $K$ boven die van buizen met de gemiddelde grootte.
\end{solution}

\begin{solution}{exo:b3:hydrogeology:11}
(a)~De logaritme verdwijnt wanneer $2.25\,Tt = r^2S$, dat wil zeggen bij $r_0 =
\sqrt{2.25\,Tt/S}$. (b)~Gespannen en freatisch:
\[
r_0 = \sqrt{\frac{2.25\times450\times30}{\num{2e-4}}} = \qty{12.3}{km},
\qquad r_0 = \sqrt{\frac{2.25\times450\times30}{0.2}} = \qty{390}{m},
\]
$\sqrt{1000} = 31.6$ keer kleiner. (c)~De druk
verspreidt zich met de diffusiviteit $T/S$, enorm wanneer $S$ klein is: binnen
de test bereikt de \omterm{def:b3:hydrogeology:cone}{depressiekegel} breuken, ontsluitingen van de watervoerende laag en rivieren,
en de verlaging wijkt af van de rechte lijn, wat de
hydrogeoloog laat zien dat een begrenzing is bereikt.
\end{solution}

\begin{solution}{exo:b3:hydrogeology:12}
(a)~Vanuit het met het water meebewegende referentiestelsel gezien is het front gespreid met
standaardafwijking $\sigma_x = \sqrt{2Dt} = \sqrt{2\alpha_Lvt}$. Het passeert
$x$ rond $t = x/v$, wanneer $\sigma_x = \sqrt{2\alpha_Lx}$; bij de snelheid
$v$, een spreiding van $\sigma_x$ in afstand komt overeen met $\sigma_t = \sigma_x/v$
in tijd. (b)~De \qty{16}{\percent} en \qty{84}{\percent} tijden liggen één
standaardafwijking aan weerszijden van het gemiddelde: $\sigma_t = (247.2 -
152.6)/2 = \num{47.3}\text{ dagen}$. Dan $\alpha_L = v^2\sigma_t^2/(2x) =
0.25\times47.3^2/200 = \qty{2.8}{m}$. (De kromme is berekend met
\qty{3}{m}: de normaalbenadering onderschat hem licht.) (c)~$R =
1 + 1.9\times0.32/0.30 = \num{3.03}$: de opgeloste stof reist met een derde van de
snelheid van het water.
\end{solution}

\begin{solution}{pb:b3:hydrogeology:1}
\textbf{1.} Schalie $363.2 - 282.5 = \qty{80.7}{m}$; kalksteen $282.5 -
102.5 = \qty{180}{m}$.

\textbf{2.} $h = 363.2 - 45.0 = \qty{318.2}{m}$; drukhoogte bovenaan
$318.2 - 282.5 = \qty{35.7}{m}$. Het water stijgt \qty{35.7}{m} boven de
bovenkant van de kalksteen, tegengehouden door de schalie: de watervoerende laag is gespannen.

\textbf{3.} Uit $h_2 - h_1 = 550a - 100b = -1.62$ en $h_3 - h_1 = 250a
- 600b = -1.44$: $a = \num{-2.715e-3}$ en $b = \num{1.269e-3}$ per
meter.

\textbf{4.} $|\nabla h| = \sqrt{a^2 + b^2} = \num{0.0030}$. Het water
stroomt langs $-\nabla h = (\num{2.715e-3}, \num{-1.269e-3})$, naar het oosten en
zuiden: azimut $\ang{180} - \arctan(2.715/1.269) = \ang{115}$.

\textbf{5.} De punten vóór ongeveer \qty{10}{min} liggen boven de lijn: daar
geldt $u > 0.05$ en ligt de Theis-kromme boven zijn logaritmische benadering
($W > -\gamma - \ln u$, zoals de figuur van de putfunctie laat zien).

\textbf{6.} De beste lijn door de punten vanaf \qty{15}{min} gaat
door \qty{1.75}{m} bij \qty{100}{min} en \qty{2.63}{m} bij \qty{1000}{min}:
$\Delta s = \qty{0.87}{m}$ per logaritmische cyclus.

\textbf{7.} $T = 2.3\times2160/(4\pi\times0.872) = \qty{453}{m^2/d}$.

\textbf{8.} De lijn daalt \qty{0.87}{m} per cyclus vanaf \qty{1.75}{m} bij
\qty{100}{min}: nul wordt bereikt $1.75/0.872 = 2.01$ cycli eerder,
$t_0 = \qty{0.98}{min} = \qty{6.83e-4}{d}$.

\textbf{9.} $S = 2.25\,Tt_0/r^2 = 2.25\times453\times\num{6.83e-4}/3600 =
\num{1.9e-4}$, binnen het bereik van gespannen watervoerende lagen.

\textbf{10.} Bij \qty{15}{min}, $u = 3600\times\num{1.9e-4}/(4\times453
\times0.0104) = \num{0.036}$, en daarna minder.

\textbf{11.} $K = 453/180 = \qty{2.5}{m/d}$, een gebroken kalksteen.

\textbf{12.} $u = 3600\times\num{2e-4}/(4\times450\times2) =
\num{2.0e-4}$, $W = -0.5772 + 8.5172 + 0.0002 = \num{7.940}$, $s =
0.382\times7.940 = \qty{3.03}{m}$, zoals genoteerd.

\textbf{13.} $u = 10^6\times\num{2e-4}/(4\times450\times30) =
\num{0.0037}$, $W = \num{5.02}$, $s = \qty{1.92}{m}$.

\textbf{14.} $r_0 = \sqrt{2.25\times450\times30/(\num{2e-4})} =
\qty{12.3}{km}$, meer dan de breedte van het gebied. De \omterm{def:b3:hydrogeology:cone}{depressiekegel}
heeft de randen van de kalksteen bereikt, zijn ontsluitingen en de breuk, die hem
tegen minder permeabel gesteente kan zetten: de oneindige watervoerende laag van Theis
beschrijft hem niet meer, en alleen de registratie van de verlagingen zal uitwijzen of een
rand de put voedt of de stroming stopt.

\textbf{15.} $2160\times30 = \qty{64800}{m^3}$. Het komt uit opslag:
het water dat vrijkomt door de uitzetting van het water en de samendrukking van
de kalksteen, over het hele oppervlak van de kegel. Integratie van $S\,s$ over
de kegel levert $Qt$ weer op.

\textbf{16.} $u = 10^6\times0.18/(4\times450\times30) = \num{3.33}$, $W
= \num{0.0086}$, $s = \qty{3}{mm}$: de boerderij zou er niets van merken. De
onttrekking zou dan water halen uit de poriën dicht bij de put, in plaats van uit de
druk van de hele watervoerende laag.

\textbf{17.} $q = 2.5\times0.0030 = \qty{0.0075}{m/d}$; $v = 0.0075/0.18
= \qty{0.042}{m/d}$, \qty{15}{m} per jaar.

\textbf{18.} $200/0.0417 = \num{4800}\text{ dagen}$, \num{13.1}~jaar; het
verontreinigde stof, \num{32.9}~jaar.

\textbf{19.} Door een cilinder met straal $r$ stroomt het debiet $Q$ door het
oppervlak $2\pi rb$ waarvan de poriën de fractie $n_e$ vormen: $\mathrm dr/\mathrm
dt = -Q/(2\pi rbn_e)$, en integratie van $2\pi bn_er\,\mathrm dr = -Q\,
\mathrm dt$ van $r$ tot 0 geeft $t = \pi n_ebr^2/Q$. Hier geldt $\pi\times0.18
\times180\times200^2/2160 = \num{1885}\text{ dagen}$, \num{5.2}\text{ jaar} voor het
water en $2.5\times5.2 = \num{12.9}\text{ jaar}$ voor het verontreinigde stof: de put
brengt het er meer dan twee keer zo snel.

\textbf{20.} $\alpha_L = \qty{20}{m}$ en $\sqrt{2\times20\times200} =
\qty{89}{m}$. De concentratie bereikt \qty{16}{\percent} van haar
eindwaarde \qty{89}{m} vóór het gemiddelde front: dat niveau bereikt de put
wanneer het gemiddelde front nog \qty{89}{m} verwijderd is, na $(111/200)^2\times12.9 = \num{4.0}\text{ jaar}$, een derde van de
gemiddelde tijd. De eerste sporen zijn belangrijk voor drinkwater, niet het
gemiddelde.

\textbf{21.} Voor: de put vangt het verontreinigde stof op, zodat het zich niet verder
stroomafwaarts kan verspreiden, en het water kan worden behandeld waar het wordt opgepompt.
Tegen: de onttrekking trekt het verontreinigde stof eerder naar een drinkwaterput dan
het anders ergens naartoe zou bewegen, en het water moet tientallen jaar worden gemonitord en
behandeld.

\textbf{22.} De \omterm{def:b3:hydrogeology:storage}{transmissiviteit} van de kalksteen van Varden is $T \approx
\textbf{\qty[text-series-to-math=true]{450}{m^2/d}}$ (\qty{453}{m^2/d}
uit de test).
\end{solution}
